\documentclass[10pt,a4paper]{amsart}
\usepackage[margin=19mm]{geometry}
\usepackage{amsmath,amssymb,mathtools}
\usepackage[scr=rsfs,cal=euler]{mathalfa}
\usepackage{xfrac}
\usepackage[unicode,hidelinks,bookmarks=false]{hyperref}
\newcommand{\ie}{i.e.\ }
\newcommand{\cf}{cf.\ }
\newcommand{\resp}{respectively\ }
\newcommand{\Subsection}{\subsection}
\providecommand{\nobreakdash}{\nobreak\hbox{-}}
\setlength{\emergencystretch}{2em}
\multlinegap0pt
\newtheorem{theorem}{Theorem}
\title{Filters, topologies, the Rado graph, and the Urysohn space}
\author{Peter J. Cameron}
\date{Source: arXiv:2602.22864v1 (CC BY 4.0)}
\begin{document}
\maketitle

\noindent\emph{Source and licence.} Filters, topologies, the Rado graph, and the Urysohn space. Version arXiv:2602.22864v1 (CC BY 4.0), distributed by arXiv under the Creative Commons Attribution 4.0 International licence, http://creativecommons.org/licenses/by/4.0/, as recorded in the arXiv licence metadata for this exact version and shown on the abstract page https://arxiv.org/abs/2602.22864v1. Full licence notice: \texttt{licenses/arxiv-2602.22864-CC-BY-4.0.txt}.\par\smallskip

\noindent\textit{Editorial scope.} Counted unit: the article's theorem theorem (source0.tex lines 271--278). The article's complete original proof of it is delivered with it (source0.tex lines 280--289, one \texttt{proof} environment of 10 lines). No mathematics is written, repaired or re-derived: the statement and the proof are the article's own lines, in the article's own notation, and no retained line is substituted or re-flowed.  No further source-local context is needed: every cross-reference of every retained line resolves inside the two delivered spans.  Nothing was omitted from any retained span, so the package carries no omission record.  No retained line contains a citation, so the wrapper carries no bibliography and no bibliography entry.  No retained line contains a figure environment, an image-inclusion command or drawing code, so no image is delivered and the package declares no asset dependency.  Editorial formatting only, in addition: an AMS \texttt{amsart} preamble that restates the article's own declarations and macros used by the retained lines, namely the environments \texttt{theorem} on separate counters, together with the standard packages those lines need.  The retained environments print in this document as Theorem 1 in order of appearance, whereas the article prints the same items with the numbering of its own full document.  The complete native source of the article as distributed in the arXiv e-print for this exact version is bundled unchanged as \texttt{sources/195356/source0.tex}.
\begin{theorem}
The following three conditions on a countable graph $\Gamma$ are equivalent:
\begin{enumerate}
\item $\mathcal{F}_\Gamma$ is nontrivial;
\item $\Gamma$ contains $R$ as a spanning subgraph;
\item $\mathcal{F}_\Gamma\subseteq\mathcal{F}_R$.
\end{enumerate}
\end{theorem}

\begin{proof}
The filter $\mathcal{F}_\Gamma$ is non-trivial if and only if any finitely
many neighbourhoods in $\Gamma$ have non-empty intersection. Now we can
define an embedding of $R$ into $\Gamma$ by modifying the standard
back-and-forth proof by going forth only. Thus (a) and (b) are equivalent.

If $\Gamma$ contains $R$ as a spanning subgraph, then $R(v)\subseteq\Gamma(v)$,
so (c) holds. Conversely, the proof that (a) and (b) are equivalent and 
standard properties of $R$ show that $\mathcal{F}_R$ is non-trivial.
\end{proof}

\end{document}
